\chapter{Banques centrales et taux court}\label{ch:m2:central-banks-and-the-short-rate}

À 14:00, heure de New York, le 16 septembre 2026, le Federal Open Market
Committee a publié un communiqué de quelques paragraphes relevant sa cible pour
le taux des fonds fédéraux d'un quart de point, dans une fourchette de 3.75 à 4~\%.
Personne à la Federal Reserve n'a acheté ni vendu une obligation pour que cela soit. Dès le
lendemain matin, la Fed versait 3.90~\% sur chaque dollar que les banques conservaient
sur leurs comptes auprès d'elle, prêtait contre collatéral en titres du Trésor à 4.00~\% à
toute contrepartie éligible qui le demandait, et reprenait des liquidités aux fonds monétaires à 3.75~\%~: trois taux
qu'elle annonçait simplement, et entre lesquels évoluaient les taux de marché de la
première devise mondiale. Le même jour, la hausse d'un quart
de point de la Banque centrale européenne entrait en vigueur~; deux jours plus tard, la Banque du Japon relevait sa cible de
environ 1~\% à environ 1.25~\%. Chaque
taux d'intérêt de ce livre, le rendement d'une obligation trentenaire, le prix d'un
swap, les points de terme d'une devise, repose sur le taux au jour le jour
que fixent ces quelques institutions. Ce chapitre explique comment elles le fixent, et
comment le marché en a fait l'indice de référence qui a remplacé les \omterm{def:m2:central-banks-and-the-short-rate:ibor}{taux
interbancaires offerts} des quarante années précédentes.

\section{Réserves et bilan de la banque centrale}

\begin{definition}[Réserves]\label{def:m2:central-banks-and-the-short-rate:reserves}
Les \emph{réserves}\index{réserves} sont les soldes que détiennent les banques, ainsi que certaines autres
institutions admises auprès d'elle, sur des comptes ouverts à la banque centrale. Il s'agit
de l'actif dans lequel les paiements entre banques sont définitivement réglés~: quand un
client d'une banque règle un client d'une autre, les réserves passent du compte de la
première banque vers le compte de la seconde.
\end{definition}

Les \omterm{def:m2:central-banks-and-the-short-rate:reserves}{réserves} sont un passif de la banque centrale, au même titre que les billets qu'elle émet et
le compte qu'elle tient pour l'État. De l'autre côté de son bilan
figurent les actifs qu'elle a achetés ou prêtés contre garantie~: obligations d'État, titres
hypothécaires, prêts aux banques. La \cref{fig:m2:central-banks-and-the-short-rate:balance}
présente les deux faces du bilan de la Fed un mercredi.

\begin{proposition}[Seule la banque centrale modifie le total]\label{prop:m2:central-banks-and-the-short-rate:total}
Les paiements entre banques et avec leurs clients déplacent des \omterm{def:m2:central-banks-and-the-short-rate:reserves}{réserves} d'un compte à un autre, mais
ne modifient pas leur total. Celui-ci ne change que lorsqu'un poste du bilan de la
banque centrale évolue~: il augmente lorsque la banque centrale achète un actif ou
prête, et diminue lorsqu'elle vend ou qu'elle est remboursée, ou lorsqu'un autre de ses
passifs (monnaie fiduciaire, compte du Trésor, reverse repos) croît aux dépens des
\omterm{def:m2:central-banks-and-the-short-rate:reserves}{réserves}.
\end{proposition}
\begin{proof}
Le bilan de la banque centrale vérifie~: actifs $=$ \omterm{def:m2:central-banks-and-the-short-rate:reserves}{réserves} $+$ autres
passifs $+$ capital. Un paiement entre deux banques débite un compte de \omterm{def:m2:central-banks-and-the-short-rate:reserves}{réserves}
et en crédite un autre~: aucun terme de l'identité ne bouge. Les \omterm{def:m2:central-banks-and-the-short-rate:reserves}{réserves} ne peuvent
donc varier que si les actifs, le capital ou un autre passif varient.
\end{proof}

\begin{example}[Jour de paiement des impôts]\label{ex:m2:central-banks-and-the-short-rate:tax}
Les entreprises versent 100~milliards de dollars d'impôts sur le compte du Trésor à la Fed.
La banque de chaque payeur perd des \omterm{def:m2:central-banks-and-the-short-rate:reserves}{réserves}~; le compte du Trésor en gagne le
même montant. Le système bancaire dispose de 100~milliards de dollars de \omterm{def:m2:central-banks-and-the-short-rate:reserves}{réserves} en moins, alors que
personne n'a rien décidé en matière de politique monétaire. Les mouvements
du compte de l'État figurent parmi les plus importants de ces \emph{facteurs autonomes}, et une banque centrale
qui veut un taux au jour le jour stable doit soit les compenser, soit détenir
tant de \omterm{def:m2:central-banks-and-the-short-rate:reserves}{réserves} qu'elles ne comptent pas.
\end{example}

\begin{omfigure}
\begin{tikzpicture}
\begin{axis}[width=11cm,height=5.2cm,ybar stacked,bar width=34pt,
  symbolic x coords={assets,liabilities},xtick=data,xticklabels={actif,passif},
  ylabel={milliards de dollars},ymin=0,ymax=7200,enlarge x limits=0.55,
  tick label style={font=\small},label style={font=\small},
  legend columns=4,legend cell align=left,legend style={font=\footnotesize,at={(0.5,-0.2)},anchor=north,draw=none,fill=none,/tikz/every even column/.append style={column sep=6pt}},
  table/col sep=comma]
\addplot[fill=omDef!70,draw=omDef] table[x=side,y=treasuries]{figdata/markets-2/01-central-banks-and-the-short-rate/balance.csv};
\addplot[fill=omDef!35,draw=omDef] table[x=side,y=mbs]{figdata/markets-2/01-central-banks-and-the-short-rate/balance.csv};
\addplot[fill=gray!30,draw=gray] table[x=side,y=otherassets]{figdata/markets-2/01-central-banks-and-the-short-rate/balance.csv};
\addplot[fill=omThm!70,draw=omThm] table[x=side,y=reserves]{figdata/markets-2/01-central-banks-and-the-short-rate/balance.csv};
\addplot[fill=omProp!60,draw=omProp] table[x=side,y=tga]{figdata/markets-2/01-central-banks-and-the-short-rate/balance.csv};
\addplot[fill=omProp!25,draw=omProp] table[x=side,y=rrp]{figdata/markets-2/01-central-banks-and-the-short-rate/balance.csv};
\addplot[fill=omThm!25,draw=omThm] table[x=side,y=currency]{figdata/markets-2/01-central-banks-and-the-short-rate/balance.csv};
\legend{titres du Trésor,MBS,autres actifs,réserves,compte du Trésor,reverse repos,monnaie fiduciaire et reste}
\end{axis}
\end{tikzpicture}

\omcaption{Bilan de la Federal Reserve le mercredi 16 septembre 2026~:
6\,747 milliards de dollars d'actifs, dont les deux tiers en titres du Trésor, financés par
3\,014 milliards de \omterm{def:m2:central-banks-and-the-short-rate:reserves}{réserves}, 877 milliards sur le compte du Trésor, 324
milliards de reverse repos, et la monnaie fiduciaire, les autres passifs et le capital
ensemble (le solde, 2\,532 milliards). Données~: Federal Reserve, H.4.1
publication du 17 septembre 2026.}\label{fig:m2:central-banks-and-the-short-rate:balance}
\end{omfigure}

\section{Systèmes de corridor et de plancher}

\begin{definition}[Taux directeur et facilités permanentes]\label{def:m2:central-banks-and-the-short-rate:policy}
Le \emph{taux directeur}\index{taux directeur} d'une banque centrale est le taux d'intérêt
qu'elle annonce comme orientation de sa politique monétaire, et vers laquelle elle guide les
taux au jour le jour du marché. Une \emph{facilité permanente}\index{facilité permanente}
est une opération accessible à l'initiative de ses contreparties, pour un montant
quelconque (contre collatéral lorsqu'elle prête), à un taux que fixe la banque centrale~: une
\emph{facilité de dépôt} qui rémunère les \omterm{def:m2:central-banks-and-the-short-rate:reserves}{réserves} et une \emph{facilité de
prêt} qui prête des \omterm{def:m2:central-banks-and-the-short-rate:reserves}{réserves} au jour le jour.
\end{definition}

\begin{proposition}[Les facilités encadrent le taux au jour le jour]\label{prop:m2:central-banks-and-the-short-rate:bounds}
Soient $i_D < i_L$ les taux des facilités de dépôt et de prêt. Une banque
qui a accès aux deux ne prête jamais de \omterm{def:m2:central-banks-and-the-short-rate:reserves}{réserves} au jour le jour en dessous de $i_D$ et
n'emprunte jamais au-dessus de $i_L$ (contre le même collatéral). Parmi de telles banques, le
taux au jour le jour se situe dans $[i_D, i_L]$.
\end{proposition}
\begin{proof}
Prêter sur le marché à $r < i_D$ rapporte moins que laisser les \omterm{def:m2:central-banks-and-the-short-rate:reserves}{réserves} à la
facilité de dépôt~; emprunter à $r > i_L$ coûte plus cher que la facilité de
prêt. Aucun des deux côtés n'accepte un tel échange.
\end{proof}

\begin{definition}[Systèmes de corridor et de plancher]\label{def:m2:central-banks-and-the-short-rate:corridor}
Dans un \emph{système de corridor}\index{système de corridor}, la banque centrale maintient
des \omterm{def:m2:central-banks-and-the-short-rate:reserves}{réserves} rares, de sorte que les banques les échangent entre elles et que le taux au jour le jour
se fixe dans $[i_D, i_L]$, près d'une cible qu'elle pilote en ajustant
l'offre de \omterm{def:m2:central-banks-and-the-short-rate:reserves}{réserves}, jour après jour. Dans un \emph{système de plancher}\index{système de plancher}
elle fournit des \omterm{def:m2:central-banks-and-the-short-rate:reserves}{réserves} au-delà de ce dont les banques ont besoin, de sorte que la réserve marginale
ne vaut que le taux de dépôt, et le taux au jour le jour se situe à, ou juste en dessous de,
$i_D$.
\end{definition}

La \cref{fig:m2:central-banks-and-the-short-rate:corridor} résume toute la théorie en
une seule courbe. La demande de \omterm{def:m2:central-banks-and-the-short-rate:reserves}{réserves} des banques baisse avec leur prix~: quand les \omterm{def:m2:central-banks-and-the-short-rate:reserves}{réserves}
sont rares, une banque en déficit en fin de journée paie jusqu'au taux de prêt~;
quand elles abondent, aucune banque ne paie plus que le taux de dépôt pour en détenir une
de plus. Le point où la droite d'offre coupe la courbe donne le taux au jour le jour. Sur la
partie pentue, une faible variation de l'offre fait bouger le taux, et la banque centrale
doit prévoir les facteurs autonomes de l'\cref{ex:m2:central-banks-and-the-short-rate:tax}
chaque matin. Sur la partie plate, elle peut les ignorer~: le taux est celui qu'elle
verse sur les dépôts, et déplacer ce taux déplace le marché.

\begin{omfigure}
\begin{tikzpicture}
\begin{axis}[width=11cm,height=5cm,
  xlabel={réserves fournies (unités arbitraires)},ylabel={taux au jour le jour (\%)},
  tick label style={font=\small},label style={font=\small},
  xmin=0,xmax=6,ymin=2.4,ymax=3.0,grid=major,grid style={gray!25},table/col sep=comma]
\addplot[omDef,very thick] table[x=reserves,y=rate]{figdata/markets-2/01-central-banks-and-the-short-rate/corridor.csv};
\draw[omThm,thick,dashed] (axis cs:0,2.90) -- (axis cs:6,2.90);
\draw[omThm,thick,dashed] (axis cs:0,2.50) -- (axis cs:6,2.50);
\draw[omProp,thick] (axis cs:2.6,2.4) -- (axis cs:2.6,3.0);
\draw[omProp,thick] (axis cs:5.2,2.4) -- (axis cs:5.2,3.0);
\node[font=\footnotesize,anchor=south west] at (axis cs:0.1,2.905) {facilité de prêt 2.90};
\node[font=\footnotesize,anchor=north west] at (axis cs:0.1,2.49) {facilité de dépôt 2.50};
\node[font=\footnotesize,anchor=west,text=omProp] at (axis cs:2.65,2.78) {corridor};
\node[font=\footnotesize,anchor=east,text=omProp] at (axis cs:5.15,2.78) {plancher};
\end{axis}
\end{tikzpicture}

\omcaption{Courbe de demande stylisée des \omterm{def:m2:central-banks-and-the-short-rate:reserves}{réserves}, entre les taux des facilités de la BCE
du 16 septembre 2026. Avec une offre située sur la droite verticale de gauche (\omterm{def:m2:central-banks-and-the-short-rate:reserves}{réserves}
rares), le taux se fixe sur la partie pentue et varie à chaque changement de
l'offre~; à droite (\omterm{def:m2:central-banks-and-the-short-rate:reserves}{réserves} abondantes), il repose sur le taux de dépôt. La
forme de la courbe est illustrative~; les deux taux des facilités sont ceux de la BCE.
Données~: simulation du chapitre.}\label{fig:m2:central-banks-and-the-short-rate:corridor}
\end{omfigure}

\begin{example}[Le corridor de l'euro]\label{ex:m2:central-banks-and-the-short-rate:ecb}
À partir du 16 septembre 2026, la facilité de dépôt de la BCE rémunère 2.50~\%, ses
opérations principales de refinancement hebdomadaires prêtent à 2.65~\% et sa facilité de prêt marginale
à 2.90~\%. Le corridor a une largeur de 40 points de base. Depuis le 18 septembre 2024, le
taux de refinancement ne se situe plus qu'à 15 points de base au-dessus du taux de dépôt (il
était de 50), une modification annoncée en mars 2024 pour que les banques empruntent lors des
opérations hebdomadaires à mesure que les \omterm{def:m2:central-banks-and-the-short-rate:reserves}{réserves} se réduisent et que le taux au jour le jour reste proche
du taux de dépôt. La zone euro a, de fait, choisi un plancher tout en
conservant le corridor comme bornes extérieures.
\end{example}

\begin{remark}[Sous le plancher]\label{rem:m2:central-banks-and-the-short-rate:below}
La borne de la \cref{prop:m2:central-banks-and-the-short-rate:bounds} ne tient que
pour les établissements qui peuvent utiliser la facilité de dépôt. Les établissements riches en liquidités
qui ne le peuvent pas, comme les fonds monétaires, prêtent aux banques au taux que
celles-ci consentent à payer, et les banques déposent les sommes reçues auprès de la banque centrale pour la
différence. C'est pourquoi un \omterm{def:m2:central-banks-and-the-short-rate:rfr}{taux de référence au jour le jour} peut s'établir juste sous un plancher~:
le 22 septembre 2026, le taux à court terme de la zone euro s'élevait à 2.440~\%, soit six points de base
sous le taux de dépôt. La réponse de la Fed à cette même fuite est sa facilité de reverse repo,
ouverte aux fonds monétaires et aux entreprises parrainées par l'État, ainsi
que les banques et les dealers, ce qui place un second plancher, plus bas, sous le marché.
\end{remark}

\begin{dated}{2026-09}{Taux administrés de quatre banques centrales}\label{dat:m2:central-banks-and-the-short-rate:rates}
\textbf{Federal Reserve} (à partir du 17 septembre 2026)~: fourchette cible pour le
taux des fonds fédéraux 3.75--4.00~\%~; intérêt servi sur les soldes de \omterm{def:m2:central-banks-and-the-short-rate:reserves}{réserves} 3.90~\%~;
\omterm{def:m2:central-banks-and-the-short-rate:policy}{facilité permanente} de repo 4.00~\%~; facilité de reverse repo au jour le jour 3.75~\% (160
milliards de dollars par contrepartie et par jour)~; crédit primaire 4.00~\%.
\textbf{BCE} (à partir du 16 septembre 2026)~: facilité de dépôt 2.50~\%, opérations principales de
refinancement 2.65~\%, facilité de prêt marginale 2.90~\%.
\textbf{Bank of England}~: Bank Rate 3.75~\%, maintenu le 16 septembre 2026 par six
voix contre trois. \textbf{Bank of Japan} (à partir du 24 septembre 2026)~: cible pour
taux monétaire non garanti au jour le jour, autour de 1.25~\%~; facilité de dépôt
complémentaire 1.25~\%~; taux de prêt de base 1.50~\%.
\end{dated}

\section{Opérations d'open market et politique de bilan}

\begin{definition}[Opération d'open market]\label{def:m2:central-banks-and-the-short-rate:omo}
Une \emph{opération d'open market}\index{opération d'open market} est un achat ou
une vente de titres, ou une opération de repo ou de reverse repo, que la banque centrale
initie, pour un montant qu'elle choisit, afin de modifier l'offre de \omterm{def:m2:central-banks-and-the-short-rate:reserves}{réserves}.
\end{definition}

\begin{definition}[Assouplissement quantitatif]\label{def:m2:central-banks-and-the-short-rate:qe}
L'\emph{assouplissement quantitatif}\index{assouplissement quantitatif} est l'achat de
montants importants de titres à plus long terme, d'obligations d'État et parfois
d'obligations hypothécaires ou d'obligations d'entreprises, afin de faire baisser les rendements à plus long terme une fois que le \omterm{def:m2:central-banks-and-the-short-rate:policy}{taux directeur}
ne peut plus baisser, ou ne peut plus baisser utilement. Son inverse, qui consiste à laisser les titres
arriver à échéance sans les remplacer ni les vendre, est le \emph{resserrement
quantitatif}.
\end{definition}

\begin{example}[Un achat, trois bilans]\label{ex:m2:central-banks-and-the-short-rate:purchase}
La banque centrale achète des obligations pour 10~milliards de dollars auprès d'un fonds de pension. La
banque centrale~: titres $+10$, \omterm{def:m2:central-banks-and-the-short-rate:reserves}{réserves} (de la banque du fonds) $+10$. La
banque du fonds~: \omterm{def:m2:central-banks-and-the-short-rate:reserves}{réserves} $+10$, dépôt du fonds $+10$. Le fonds~: obligations
$-10$, dépôt $+10$. L'achat a créé des \omterm{def:m2:central-banks-and-the-short-rate:reserves}{réserves} et un dépôt bancaire de
même montant~; il n'a pas créé de prêt, et la banque ne peut pas «~prêter~» les
\omterm{def:m2:central-banks-and-the-short-rate:reserves}{réserves} à quiconque sauf à une autre banque, où elles restent des \omterm{def:m2:central-banks-and-the-short-rate:reserves}{réserves}
(\cref{prop:m2:central-banks-and-the-short-rate:total}).
\end{example}

Dans un \omterm{def:m2:central-banks-and-the-short-rate:corridor}{système de corridor}, les \omterm{def:m2:central-banks-and-the-short-rate:omo}{opérations d'open market} \emph{constituent} la politique monétaire~: la
banque centrale ajoute ou retire des \omterm{def:m2:central-banks-and-the-short-rate:reserves}{réserves} chaque jour pour maintenir le taux à sa cible. Dans un
\omterm{def:m2:central-banks-and-the-short-rate:corridor}{système de plancher}, elles ne le sont pas~: le taux est fixé par le taux de dépôt, et la taille
du bilan devient un instrument distinct, utilisé pour les rendements à plus long terme
(voir la \cref{def:m2:central-banks-and-the-short-rate:qe}) ou pour conserver
des \omterm{def:m2:central-banks-and-the-short-rate:reserves}{réserves} en toute sécurité sur la partie plate de la courbe. Cette séparation est le
principe d'un plancher. Elle en est aussi le coût~: une banque centrale qui verse des intérêts sur
des milliers de milliards de \omterm{def:m2:central-banks-and-the-short-rate:reserves}{réserves} rémunère les banques pour les détenir, au grand jour.

L'histoire récente de la Fed montre les deux instruments à l'œuvre. Elle
a réduit ses avoirs en titres de plus de 2.2~billions de dollars à partir
de juin 2022, s'est arrêtée le 1\up{er} décembre 2025 lorsqu'elle a jugé les
\omterm{def:m2:central-banks-and-the-short-rate:reserves}{réserves} proches de la limite de l'abondance, et a recommencé à acheter
des bons du Trésor en décembre 2025, non pour assouplir sa politique, mais
pour que les \omterm{def:m2:central-banks-and-the-short-rate:reserves}{réserves} continuent de croître avec l'économie. Sa note de mise
en œuvre de septembre 2026, celle-là même qui a relevé les taux, répète la
consigne d'acheter des bons selon les besoins~: le taux a monté et le bilan a continué de croître, sans contradiction.

\section{Taux de référence au jour le jour et fin des taux interbancaires offerts}

\begin{definition}[Taux interbancaire offert]\label{def:m2:central-banks-and-the-short-rate:ibor}
Un \emph{taux interbancaire offert}\index{taux interbancaire offert} (IBOR) est un
indice de référence pour les emprunts à terme non garantis entre banques, dont
l'échéance va d'un à douze mois, calculé chaque jour à partir des soumissions d'un panel de banques sur les taux
auxquels elles pouvaient emprunter. Le LIBOR, pour cinq devises, était le plus important~; des milliers de milliards
de prêts, d'obligations et de dérivés le référençaient.
\end{definition}

\begin{definition}[Taux de référence au jour le jour]\label{def:m2:central-banks-and-the-short-rate:rfr}
Un \emph{taux de référence au jour le jour}\index{taux de référence au jour le jour} (également
appelé taux sans risque, RFR) est un indice de référence calculé à partir
des transactions réelles au jour le jour, publié par une banque centrale ou un administrateur le
jour ouvré suivant. Il est soit garanti (repo sur obligations d'État), soit non garanti
(dépôts collectés par les banques).
\end{definition}

La transition du premier type vers le second a été imposée par la rareté
du marché que le premier mesurait~: peu de banques se prêtaient encore entre elles
sans garantie pendant trois mois, et un indice fondé sur le jugement plutôt que sur les
transactions s'est révélé à la fois fragile et manipulable. Le 5 mars 2021, le régulateur britannique
a annoncé les dates de fin~: la plupart des paramètres du LIBOR après le 31 décembre
2021, les principaux paramètres en dollars après le 30 juin 2023. Des paramètres synthétiques en dollars
ont été publiés en guise de pont pour les contrats existants jusqu'au 30 septembre 2024~;
depuis lors, les trente-cinq paramètres du LIBOR ont cessé d'être publiés.

\begin{dated}{2026-09}{Les quatre principaux taux de référence au jour le jour}\label{dat:m2:central-banks-and-the-short-rate:benchmarks}
\textbf{SOFR} (dollar américain, garanti)~: repo au jour le jour sur titres du Trésor, tripartite,
collatéral général et bilatéral compensé, les specials étant écartés~;
publié par la Federal Reserve Bank of New York vers 08:00, heure de New York,
pour le jour ouvré précédent. \textbf{\texteuro STR} (euro,
non garanti)~: emprunts au jour le jour des banques de la zone euro sur le marché de gros~; publiés par
la BCE chaque jour ouvré TARGET pour le précédent (22 septembre 2026~: 2.440~\%
sur 69.9~milliards d'euros, 929 transactions, 47 banques). \textbf{SONIA} (livre sterling,
non garanti)~: ce que paient les banques pour emprunter des livres sterling au jour le jour
auprès des institutions financières~; publié par la Banque d'Angleterre à 09:00, heure de Londres, le
jour ouvré suivant.
\textbf{Japon}~: le taux du call non garanti au jour le jour, la cible
opérationnelle de la Banque du Japon.
\end{dated}

Un taux au jour le jour n'est connu qu'un jour à la fois, mais un prêt ou un swap
verse des intérêts pendant un trimestre ou un an. La réponse du marché consiste à composer les taux
quotidiens sur la période et à payer à son terme.

\begin{definition}[Composé a posteriori]\label{def:m2:central-banks-and-the-short-rate:arrears}
Sur une période d'intérêts de $N$ jours calendaires, avec les jours ouvrés $d_1, \dots,
d_M$, les taux fixés $r_1, \dots, r_M$ et les poids $n_i$, les jours calendaires allant de
$d_i$ jusqu'au jour ouvré suivant, le taux \emph{composé a
posteriori}\index{composé a posteriori} (compounded in arrears) est
\[
R \;=\; \Bigl[\,\prod_{i=1}^{M}\Bigl(1 + \frac{r_i\,n_i}{B}\Bigr) - 1\Bigr]\frac{B}{N},
\]
avec $B = 360$ pour le dollar et l'euro, et 365 pour la livre sterling. Il n'est connu
qu'après le dernier taux fixé. Trois conventions avancent la date de paiement~: un
\emph{lookback} de $p$ jours ouvrés utilise $r$ observé $p$ jours avant chaque $d_i$
avec les poids propres à la période~; un \emph{décalage d'observation} (observation shift) prend les taux et
les poids d'une période décalée de $p$ jours ouvrés~; un \emph{verrouillage} (lockout)
gèle les $k$ derniers taux fixés à la valeur de celui qui les précède.
\end{definition}

\begin{proposition}[La composition ajoute un terme du second ordre]\label{prop:m2:central-banks-and-the-short-rate:convexity}
Soit $R_s = \sum_i r_i n_i / N$ la moyenne simple pondérée par les jours. Si tous les
$r_i \ge 0$, alors $R \ge R_s$~; pour un taux constant $r$ fixé chaque
jour calendaire,
\[
R - r \;=\; \frac{r^2 (N-1)}{2B} + O\!\left(\frac{r^3N^2}{B^2}\right).
\]
\end{proposition}
\begin{proof}
Avec $x_i = r_i n_i / B \ge 0$, $\prod (1 + x_i) - 1 = \sum x_i + \sum_{i<j}
x_i x_j + \dots \ge \sum x_i$, et $\sum x_i \cdot B/N = R_s$. Pour un taux
quotidien constant, $(1 + r/B)^N - 1 = Nr/B + \binom{N}{2}(r/B)^2 + O((Nr/B)^3)$~;
multiplier par $B/N$.
\end{proof}

\begin{example}[Septembre 2026, illustré]\label{ex:m2:central-banks-and-the-short-rate:month}
On prend des taux fixés illustratifs de 3.62~\% jusqu'au 16 septembre, 3.87~\% à partir du
relèvement, et 3.95~\% le jour de fin de trimestre, le 30 septembre
(\cref{fig:m2:central-banks-and-the-short-rate:month}). Du 1\up{er} septembre au
1\up{er} octobre (30 jours, 21 taux fixés, fête du Travail le 7 septembre), la moyenne pondérée par les jours
est de 3.7393~\% et le taux composé de 3.7448~\%~: la composition ajoute 0.54 points
de base, contre les 0.56 prévus par la \cref{prop:m2:central-banks-and-the-short-rate:convexity}
pour un taux plat de 3.74~\%. Avec un lookback de cinq jours, le taux vaut 3.6836~\%, inférieur de 6.12 points
de base, soit 5\,098~dollars sur 100~millions de dollars~: le lookback reporte cinq jours
du nouveau taux dans la période suivante, où l'emprunteur les paiera.
\end{example}

\begin{omfigure}
\begin{tikzpicture}
\begin{axis}[width=11cm,height=5cm,
  xlabel={jour de septembre 2026},ylabel={taux (\%)},
  tick label style={font=\small},label style={font=\small},
  xmin=1,xmax=31,ymin=3.55,ymax=4.0,grid=major,grid style={gray!25},
  legend columns=2,legend style={font=\footnotesize,at={(0.5,-0.3)},anchor=north,draw=none,fill=none,/tikz/every even column/.append style={column sep=8pt}},
  table/col sep=comma]
\addplot[omDef,very thick,const plot] table[x=day,y=fixing]{figdata/markets-2/01-central-banks-and-the-short-rate/month.csv};
\addplot[omThm,very thick,const plot,dashed] table[x=day,y=compounded]{figdata/markets-2/01-central-banks-and-the-short-rate/month.csv};
\legend{fixing du jour (reporté les jours non ouvrés),composé depuis le 1\up{er} septembre}
\end{axis}
\end{tikzpicture}

\omcaption{Les fixations illustratives de l'\cref{ex:m2:central-banks-and-the-short-rate:month}
et le taux composé depuis le début de la période jusqu'à chaque jour. Le
taux composé n'intègre une variation que progressivement~: chaque jour au
nouveau niveau ajoute un poids à une moyenne portant sur tous les jours
écoulés. Les fixations du vendredi comptent pour trois jours, celle du
vendredi précédant la fête du Travail pour quatre. Les fixations sont
illustratives, non publiées au taux SOFR. Données~: le tutoriel du chapitre.}\label{fig:m2:central-banks-and-the-short-rate:month}
\end{omfigure}

\begin{definition}[Spread de repli]\label{def:m2:central-banks-and-the-short-rate:fallback}
Un \emph{spread de repli}\index{spread de repli} est l'ajustement fixe ajouté
à un taux overnight composé lorsqu'il remplace un \omterm{def:m2:central-banks-and-the-short-rate:ibor}{taux interbancaire offert}
dans un contrat existant, afin de compenser la prime de crédit et de terme
que contenait l'ancien taux. Pour le LIBOR, il a été fixé comme la médiane sur
cinq ans de l'écart entre les deux taux, le 5 mars 2021~: pour le
LIBOR en dollars face au SOFR, 11.448 points de base à un mois, 26.161 à trois
mois et 42.826 à six mois.
\end{definition}

\begin{omfigure}
\begin{tikzpicture}[font=\small]
\draw[thick,-{Stealth[length=2.5mm]}] (0,0) -- (13,0) node[right]{temps};
\draw[omDef,line width=4pt] (3.5,0.55) -- (11.5,0.55);
\node[above=3pt,text=omDef,font=\footnotesize] at (7.5,0.6) {période d'intérêt~: poids $n_i$};
\draw[omThm,line width=4pt] (1.5,-0.55) -- (9.5,-0.55);
\node[below=3pt,text=omThm,font=\footnotesize] at (5.5,-0.6) {période d'observation~: taux $r_i$ (et, avec un décalage, poids)};
\foreach \x/\t in {3.5/début,11.5/fin}{\draw[thick] (\x,0.15) -- (\x,-0.15); \node[anchor=south,font=\footnotesize\bfseries] at (\x,0.8) {\t};}
\draw[-{Stealth[length=2mm]},omProp,thick] (3.5,-0.25) -- (1.6,-0.25);
\draw[-{Stealth[length=2mm]},omProp,thick] (11.5,-0.25) -- (9.6,-0.25);
\node[font=\footnotesize,text=omProp] at (10.5,-1.45) {$p$ jours ouvrés};
\draw[decorate,decoration={brace,amplitude=4pt},thick] (9.5,1.55) -- (11.5,1.55) node[midway,above=5pt,font=\footnotesize,align=center] {dernier taux connu~: paiement\\calculé, envoyé à la date de fin};
\end{tikzpicture}

\omcaption{Un lookback~: chaque jour de la période d'intérêt reprend la
fixation de $p$ jours ouvrés plus tôt, de sorte que le dernier taux nécessaire
est connu $p$ jours avant la fin de la période et que le paiement peut être
calculé et envoyé à sa date de fin. Avec un décalage d'observation, les
poids eux aussi proviennent de la période inférieure.}\label{fig:m2:central-banks-and-the-short-rate:lookback}
\end{omfigure}

\section{Tutoriel~: capitalisation d'un mois a posteriori}\label{tut:m2:central-banks-and-the-short-rate:compound}

\begin{tutorial}
\textbf{Objectif.} Capitaliser un mois de fixations overnight sur un calendrier
réel, avec et sans lookback, et comparer à la moyenne simple.
\textbf{État final~:} les quatre nombres de
l'\cref{ex:m2:central-banks-and-the-short-rate:month} et de
la \cref{fig:m2:central-banks-and-the-short-rate:month}.
\begin{tutsteps}
\item \textbf{Jours ouvrés et poids.} Une fixation court jusqu'au jour ouvré
suivant~; celle d'un vendredi court sur trois jours.
\omcode{code/firm/rfr/firm_rfr.py}{32}{45}{Jours ouvrés d'une période et jours calendaires couverts par chaque fixation.}\label{lst:m2:central-banks-and-the-short-rate:weights}
On doit obtenir les poids \texttt{[1, 4, 1, 1]} pour la période du 3 au 10 septembre 2026~: le
vendredi précédant la fête du Travail compte quatre jours.
\item \textbf{Quelle fixation, quel poids.} Une fonction renvoie, pour chaque
étape de la période, la date de la fixation utilisée et les jours qu'elle
couvre, pour les trois conventions.
\omcode{code/firm/rfr/firm_rfr.py}{59}{72}{Dates de fixation et poids avec un lookback, un décalage d'observation ou un lockout.}\label{lst:m2:central-banks-and-the-short-rate:obs}
\item \textbf{Capitaliser.}
\omcode{code/firm/rfr/firm_rfr.py}{75}{82}{Taux composé de la période, annualisé.}\label{lst:m2:central-banks-and-the-short-rate:compound}
\item \textbf{Lancer le mois.} \texttt{short\_rate\_demo.month\_summary()}
affiche 3.7393~\% (simple), 3.7448~\% (composé) et 3.6836~\% (lookback de
cinq jours)~; \texttt{fig\_short\_rate.py} écrit les données du graphique.
\end{tutsteps}
\textbf{Pour aller plus loin.} Remplacer le lookback par un lockout de deux
jours, et par un décalage d'observation de cinq jours, puis expliquer pourquoi
le second donne exactement le nombre du lookback sur ce mois
(\cref{exo:m2:central-banks-and-the-short-rate:7}). Ensuite, déplacer la hausse à la dernière semaine et observer l'erreur du lookback.
\end{tutorial}

\section{Construction~: le moteur de capitalisation du taux overnight}\label{bld:m2:central-banks-and-the-short-rate:rfr}

\begin{build}
\bfield{Rôle} Chaque jambe variable de la petite entreprise, prêts, swaps
(\cref{ch:m2:interest-rate-swaps}), contrats à terme sur taux overnight
(\cref{ch:m2:short-term-interest-rate-futures}), intérêts courus de repo, paie un
taux overnight composé. Ce composant le calcule, une fois pour toutes et correctement.
\bfield{Interface} \texttt{Calendar(holidays)} avec \texttt{is\_business\_day},
\texttt{add(d, n)}, \texttt{business\_days(start, end)}~;
\texttt{Convention(lookback, shift, lockout, basis)}~;
\texttt{observations(cal, start, end, conv)}~;
\texttt{compound\_in\_arrears(fixings, cal, start, end, conv)}~;
\texttt{simple\_average(\ldots)}~; \texttt{rate\_from\_index(i0, i1, start,
end)}.
\bfield{Règles} Les jours ouvrés excluent les week-ends et les jours fériés du calendrier~;
les poids courent jusqu'au jour ouvré suivant et s'arrêtent à la fin de la période~; les taux sont des décimaux
et peuvent être négatifs~; la base vaut 360 sauf si la devise indique 365.
\bfield{Tests d'acceptation} \texttt{code/firm/rfr/tests/}~: un vendredi compte
trois jours et le vendredi précédant un lundi férié, quatre~; un taux constant
se capitalise selon la forme fermée et dépasse sa moyenne simple~; un lookback
conserve les poids de la période~; un décalage d'observation reprend ceux de la
période décalée~; un lockout répète la dernière fixation libre~; la voie par
l'indice publié donne le même taux~; les taux négatifs.
\bfield{Pour aller plus loin} Arrondi comme le fait un indice composé publié
(à huit décimales)~; une convention de retard de paiement~; un calendrier de
jours fériés chargé depuis les données de référence de la petite entreprise (livre~1, chapitre~28).
\end{build}

\begin{omsources}
\item Board of Governors of the Federal Reserve System, \emph{Implementation Note issued September 16, 2026}~; communiqué statistique H.4.1, 17 septembre
2026~; procès-verbaux du FOMC, décembre 2025.
\item European Central Bank, \emph{Key ECB interest rates}~; communiqué de presse \emph{Changes to the operational framework for implementing monetary
policy}, 13 mars 2024~; \emph{Euro short-term rate}.
\item Bank of England, \emph{Monetary Policy Summary}, septembre 2026~; \emph{SONIA interest rate benchmark}. Bank of Japan, \emph{Change in the
Guideline for Money Market Operations}, 18 septembre 2026.
\item Federal Reserve Bank of New York, \emph{Secured Overnight Financing
Rate}~; \emph{Reverse repo FAQ}.
\item Financial Conduct Authority, \emph{The end of LIBOR}, 2024~; ISDA,
déclaration sur l'annonce de la FCA, 5 mars 2021~; ARRC recommended spread
adjustments, 2021.
\end{omsources}

\section{Exercices}

\begin{exercise}[$\star$]\label{exo:m2:central-banks-and-the-short-rate:1}
Une période d'intérêts court du jeudi au mardi suivant, avec des fixings
de 3.90~\% (jeudi), 3.88~\% (vendredi) et 3.92~\% (lundi). Calculer le taux
\omterm{def:m2:central-banks-and-the-short-rate:arrears}{composé a posteriori} (base 360) et la moyenne simple.
\end{exercise}

\begin{exercise}[$\star$]\label{exo:m2:central-banks-and-the-short-rate:2}
Avec les taux de la BCE de l'\cref{dat:m2:central-banks-and-the-short-rate:rates},
calculer la largeur du corridor et l'écart entre le taux de refinancement
et le taux de la facilité de dépôt. L'euro short-term rate s'établit à 2.440~\%. Quelle
borne de la \cref{prop:m2:central-banks-and-the-short-rate:bounds} semble
violée, et pourquoi ne l'est-elle pas~?
\end{exercise}

\begin{exercise}[$\star$]\label{exo:m2:central-banks-and-the-short-rate:3}
La banque centrale vend à une compagnie d'assurance des obligations pour
10~milliards de dollars. Écrire les variations des bilans de la banque centrale,
de la banque de l'assureur et de l'assureur.
\end{exercise}

\begin{exercise}[$\star\star$]\label{exo:m2:central-banks-and-the-short-rate:4}
Estimer, en points de base, l'écart entre un taux overnight plat de 4~\% composé
quotidiennement sur 91~jours calendaires et sa moyenne simple, d'abord avec la
\cref{prop:m2:central-banks-and-the-short-rate:convexity}, puis exactement.
\end{exercise}

\begin{exercise}[$\star\star$]\label{exo:m2:central-banks-and-the-short-rate:5}
Des recettes fiscales de 150~milliards de dollars sont versées sur le compte
du Trésor en une seule journée. Que se passe-t-il pour le taux overnight dans
un \omterm{def:m2:central-banks-and-the-short-rate:corridor}{système de corridor} et dans un \omterm{def:m2:central-banks-and-the-short-rate:corridor}{système de plancher}~? Que doit faire chaque banque centrale pour y remédier~?
\end{exercise}

\begin{exercise}[$\star\star$]\label{exo:m2:central-banks-and-the-short-rate:6}
Un ancien prêt de 50~millions de dollars rémunéré au LIBOR dollar à trois mois
augmenté de 150 points de base a basculé sur le SOFR composé. Sur 91~jours, le
SOFR se compose à 3.80~\%. Calculer le taux tout compris et l'intérêt (ACT/360).
\end{exercise}

\begin{exercise}[$\star\star\star$]\label{exo:m2:central-banks-and-the-short-rate:7}
\emph{Programmation.} Avec \texttt{firm\_rfr} et les fixings du tutoriel,
calculer le taux de septembre avec un délai de blocage de deux jours et avec un
décalage d'observation de cinq jours. Donner chaque écart par rapport au taux
composé simple, en points de base, et expliquer pourquoi le décalage donne exactement le nombre du lookback.
\end{exercise}

\begin{exercise}[$\star\star\star$]\label{exo:m2:central-banks-and-the-short-rate:8}
\emph{Repérer l'erreur.} «~Pour relever le taux des fonds fédéraux d'un quart de
point, le desk de trading de la Fed vend des bons du Trésor chaque matin jusqu'à
ce que le taux ait suffisamment monté.~» Corriger l'énoncé pour le système que la
Fed utilisait en 2026. Dans quel système aurait-il été à peu près juste~?
\end{exercise}

\section{Problème~: Le prix d'une réserve}

\begin{problem}[{Problème du week-end --- le mois où le corridor a bougé}]\label{pb:m2:central-banks-and-the-short-rate:1}
Une banque centrale exploite un corridor, avec une facilité de dépôt à 2.25~\% et une
facilité de prêt à 2.65~\% (base 360). Au milieu de la période de 30~jours, elle
déplace les deux de 25~points de base, pour les porter à 2.50~\% et 2.90~\%. La banque A
détient 2~milliards d'euros de \omterm{def:m2:central-banks-and-the-short-rate:reserves}{réserves} de plus que nécessaire sur toute la période~; la banque B
en manque 2~milliards d'euros sur toute la période. Il n'existe aucune autre banque.

\medskip\noindent\textbf{Partie I --- Les facilités.}
\begin{enumerate}
\item La banque A laisse son excédent à la facilité de dépôt pendant 30~jours. Que
gagne-t-elle~?
\item La banque B couvre son déficit à la facilité de prêt pendant 30~jours. Que
paie-t-elle~?
\item Quelle est la différence, et à qui profite-t-elle~?
\item Pourquoi une banque centrale laisse-t-elle un écart entre les deux taux~?
\item Les deux taux ont bougé de 25~points de base. Pourquoi cela déplace-t-il le
taux du marché de 25~points de base, quel que soit le volume des \omterm{def:m2:central-banks-and-the-short-rate:reserves}{réserves}~?
\end{enumerate}

\medskip\noindent\textbf{Partie II --- Le marché interbancaire.}
\begin{enumerate}[resume]
\item La banque A prête directement son excédent à la banque B au milieu du
corridor. Quel intérêt est versé entre les deux sur la période~?
\item Donner le gain de A par rapport à la facilité de dépôt et l'économie de B
par rapport à la facilité de prêt.
\item Rien n'impose le milieu. Qu'est-ce qui détermine où, dans le corridor, le
taux se fixe~?
\item A demande 5~points de base au-dessus du milieu pour prêter sans garantie à
B. Donner les deux nouveaux gains.
\item B ne peut emprunter à la banque centrale qu'en apportant un collatéral~; A
prête sans garantie. Quel taux mesure un indice de référence overnight non garanti, et lequel
mesure un indice garanti~?
\end{enumerate}

\medskip\noindent\textbf{Partie III --- L'inondation.}
\begin{enumerate}[resume]
\item La banque centrale achète pour 50~milliards d'euros d'obligations. Où
va le taux interbancaire, et pourquoi~?
\item Quel est le gain de A à prêter désormais à B~?
\item La banque centrale veut relever les taux de 25~points de base. Que
change-t-elle, et que n'a-t-elle plus besoin de faire~?
\item A-t-elle encore besoin d'\omterm{def:m2:central-banks-and-the-short-rate:omo}{opérations d'open market} au jour le jour~? Que
représente son bilan désormais~?
\item Un fonds monétaire disposant de 1~milliard d'euros ne peut pas utiliser la
facilité de dépôt. Il prête à A à 6~points de base sous le taux de dépôt pendant
30~jours, et A dépose les liquidités. Que gagne A sur l'opération~?
\end{enumerate}

\medskip\noindent\textbf{Partie IV --- Jugement.}
\begin{enumerate}[resume]
\item Donner deux raisons pour lesquelles une banque centrale pourrait préférer un
corridor avec des \omterm{def:m2:central-banks-and-the-short-rate:reserves}{réserves} rares.
\item Donner deux raisons pour lesquelles elle pourrait préférer un plancher.
\item Que nous apprend sur son choix l'écart de 15~points de base entre ses taux
de refinancement et de dépôt depuis 2024~?
\item Énoncer le \emph{résultat nommé}~: le gain tiré de l'échange de \omterm{def:m2:central-banks-and-the-short-rate:reserves}{réserves} sur la
période, sous forme de formule et en euros, et ce qu'il devient dans un plancher.
\item En une phrase~: qu'est-ce qui fixe le taux overnight~?
\end{enumerate}
\end{problem}

\section{Questions d'entretien}

\begin{interviewq}[$\star$ \iqroles{trader, banque}]\label{iq:m2:central-banks-and-the-short-rate:1}
Quelle est la différence entre le SOFR et le LIBOR en dollars, et pourquoi
la marge d'un prêt indexé sur le SOFR est-elle supérieure à celle du prêt LIBOR qu'il a remplacé~?
\end{interviewq}

\begin{interviewq}[$\star$ \iqroles{trader, chercheur}]\label{iq:m2:central-banks-and-the-short-rate:2}
Comment un \omterm{def:m2:central-banks-and-the-short-rate:rfr}{taux de référence au jour le jour} peut-il afficher une valeur inférieure au taux de dépôt de la banque centrale~?
\end{interviewq}

\begin{interviewq}[$\star\star$ \iqroles{trader, chercheur, banque}]\label{iq:m2:central-banks-and-the-short-rate:3}
Décrivez-moi les bilans lorsqu'une banque centrale mène un
\omterm{def:m2:central-banks-and-the-short-rate:qe}{assouplissement quantitatif}. Met-elle davantage de fonds à la disposition des banques pour prêter~?
\end{interviewq}

\begin{interviewq}[$\star\star$ \iqroles{développeur, banque}]\label{iq:m2:central-banks-and-the-short-rate:4}
Pourquoi les prêts indexés sur du SOFR \omterm{def:m2:central-banks-and-the-short-rate:arrears}{composé a posteriori} recourent-ils à un
lookback ou à un décalage d'observation~? Quelle est la différence entre les deux~?
\end{interviewq}

\begin{interviewq}[$\star\star$ \iqroles{chercheur, trader}]\label{iq:m2:central-banks-and-the-short-rate:5}
À partir des seules données de marché, comment sauriez-vous distinguer si une
banque centrale applique un \omterm{def:m2:central-banks-and-the-short-rate:corridor}{système de plancher} ou un \omterm{def:m2:central-banks-and-the-short-rate:corridor}{système de corridor}~?
\end{interviewq}

\begin{interviewq}[$\star\star\star$ \iqroles{développeur}]\label{iq:m2:central-banks-and-the-short-rate:6}
Vous écrivez la fonction qui calcule les intérêts d'un prêt indexé sur du SOFR
\omterm{def:m2:central-banks-and-the-short-rate:arrears}{composé a posteriori}. Énumérez les cas limites que vos tests doivent couvrir.
\end{interviewq}
